% =====================================================================
\chapter[Internal Infrastructuralisation]{Internal Infrastructuralisation: The Conditional Necessity of the Six Functions}\label{ch:infrastructure}
% =====================================================================

This chapter addresses Problem \ref{prob:scale}. The central claim is that any sufficiently
complex software system reinvents an operating system inside its own problem domain. The
cause is not that engineers imitate a kernel; it is that a set of semantic obligations
cannot be evaded at scale. The chapter gives the premises, the proof and the criteria for
that claim, together with a falsifiable test that can be carried out inside a single
project.

\section{The Phenomenon and the Terms}\label{sec:phenomenon}

The observed phenomenon is this: once a system grows past a certain size, it begins to
promote mechanisms that were previously scattered across it into unified system-level
capabilities. Specifically:

\begin{itemize}
\item several facilities need to allocate and release resources;
\item several subsystems need to manage object lifetime;
\item several components need to execute asynchronous tasks;
\item several modules need to share state and coordinate modifications;
\item several facilities need unified error handling, logging and diagnostics;
\item several execution paths need to obey the same thread-safety and concurrency rules.
\end{itemize}

If each module solved these problems on its own, the system would accumulate a mass of
duplicated implementations, inconsistent behaviour and coupling that is hard to control. The
system therefore begins to handle them centrally, producing common mechanisms, interfaces,
managers, runtimes and execution models.

\begin{definition}[Internal infrastructuralisation]\label{def:infra}
\term{Internal infrastructuralisation} is a system's passage through three successive
stages:
\begin{enumerate}[label=(\roman*)]
\item \term{commonality identification}: recurrent requirements are identified as one class
      of problem (a merger in problem space);
\item \term{mechanism unification}: the commonality is promoted to a common capability, so
      that different facilities access shared resources through common interfaces and obey
      common lifecycle, error-handling and concurrency rules (centralisation and
      standardisation of mechanism);
\item \term{structural unification}: upper-layer facilities no longer solve all lower-level
      problems directly, but build on unified rules (unification of structure).
\end{enumerate}
In shorthand,
{\small$\boxed{\text{commonality identification} \to \text{mechanism unification} \to
\text{upper-layer composition}}$}.
\end{definition}

Mechanism unification does \emph{not} mean that all logic is gathered into one central
module. A mature architecture can achieve unification \emph{through protocols, interfaces
and distributed ownership}. What is unified is the \emph{rules and contracts}, not
necessarily the \emph{implementations and ownership}. That distinction is the key to
avoiding a god object, and Chapter \ref{ch:design} states it as a design principle.

\section{Participant--Resource Structure}\label{sec:participant}

\begin{definition}[Participant--resource structure]\label{def:pr-structure}
A \term{participant--resource structure} is a five-tuple
\begin{equation}\label{eq:pr}
\Pi = \bigl(E,\; \mathcal{R},\; \to,\; \mathrm{Ready},\; \mathrm{Fault}\bigr),
\end{equation}
where $E$ is a set of concurrently live entities,
$\mathcal{R} = \{r_1,\dots,r_m\}$ is a finite set of shared resources, $\to$ is the
transition relation of the operational semantics,
$\mathrm{Ready} \subseteq E \times \mathrm{Conf}$ indicates which entities may execute in a
given configuration, and $\mathrm{Fault}$ is a set of fault events.
\end{definition}

The two quantities that matter are $n = |E|$ and $m = |\mathcal{R}|$. Section
\ref{sec:integration} proves that the independent variable deciding whether a unified
infrastructure is needed is the structural complexity $n \times m$, not lines of code.

\begin{definition}[Unified substrate]\label{def:substrate}
A \term{unified substrate} for $\Pi$ is a set of mechanisms $\mathcal{U}$ such that the
interaction of every entity $e \in E$ with every resource $r \in \mathcal{R}$ proceeds
through an interface declared by $\mathcal{U}$ rather than being implemented directly by
$e$.
\end{definition}

\section{The Integration Cost Theorem and the Phase-Transition Threshold}\label{sec:integration}

\begin{theorem}[Integration cost]\label{th:integration}
Let the participant--resource structure $\Pi$ have $n$ entities and $m$ shared resources.
Assume:
\begin{enumerate}[label=(\alph*)]
\item without a substrate, every entity that participates in a resource $r$ must implement
      correct interaction with $r$ independently, including independent handling of the
      invariants of $r$;
\item under maximal coupling, every entity participates in all $m$ resources;
\item with a substrate, every entity uses at most $k$ substrate interfaces, and the
      substrate itself defines the semantics of each of the $m$ resources once; the fixed
      cost of the substrate is $F$ (design, implementation and long-term maintenance).
\end{enumerate}
Then the number of interacting pairs is $\Theta(nm)$ without a substrate and
$\mathcal{O}(nk + m)$ with one. The substrate is better than no substrate in expectation if
and only if
\begin{equation}\label{eq:threshold}
n \cdot m \;>\; c\,(n k + m) \;+\; F
\end{equation}
where $c$ is the unit cost coefficient of one interaction.
\end{theorem}

\begin{proof}
Without a substrate: correct interaction of entity $e$ with resource $r$ requires an
independent implementation, so the number of interacting pairs is
$\sum_{e} |\mathcal{R}(e)| \le nm$, with equality under assumption (b), that is
$\Theta(nm)$. In addition, for each resource $r$ the $n$ independent implementations must
agree on the invariants of $r$; this is a consistency obligation whose number of failure
modes is $\Omega(n)$ (any single implementation that deviates breaks the global invariant).

With a substrate: the semantics of the $m$ resources are defined once each, at cost
$\Theta(m)$; every entity interacts with the substrate through at most $k$ interfaces, at
cost $\mathcal{O}(nk)$; add the fixed cost $F$. The total is therefore
$\mathcal{O}(nk + m) + F$.

Comparing the two expressions and normalising by the unit cost coefficient $c$ gives
\eqref{eq:threshold}.
\end{proof}

\begin{corollary}[Phase-transition threshold]\label{cor:threshold}
Because $k$ is the \emph{bounded} interface usage of each entity, $\mathcal{O}(nk + m)$ is
linear in $n$, whereas $\Theta(nm)$ is a product linear in both $n$ and $m$. Hence when
either $n$ or $m$ grows substantially, inequality \eqref{eq:threshold} must hold. The phase
transition occurs at $n m \sim F/c$, that is, when the \emph{product of the number of
entities and the number of resources} exceeds the order of the fixed cost of the substrate.
\end{corollary}

Corollary \ref{cor:threshold} makes precise the intuition that ``the code volume reaches
some order of magnitude'', and supplies its correct form:

\begin{principle}[The threshold is not lines of code]\label{pr:notloc}
The independent variable deciding whether a unified substrate must be built is the product
of the \emph{number of live entities} and the \emph{number of shared resources}, together
with the \emph{frequency of recurrent problems}, the \emph{number of interactions between
state and resources}, the \emph{complexity of concurrent execution}, the \emph{rate of
module evolution}, and the \emph{demands on consistency, performance and reliability}.
Lines of code are only a proxy for these quantities.
\end{principle}

This explains a fact that ought to be puzzling: two programs of a million lines each may
have entirely different internal architectures. One may be a large body of mutually
independent generated code, data files or linear business logic ($m$ small, $n$ small); the
other may have complex shared state, asynchronous tasks and many interdependent subsystems
($n$ and $m$ both large). Only the second develops a need for a unified runtime, resource
management and coordination mechanisms.

The central reduction (Section \ref{sec:participant}) supplies a second explanation here:
when the number of live entities in a system jumps from $10^{2}$ to $10^{6}$, naive
management undergoes a \emph{phase transition}. For 100 objects, conditional branching and
raw allocation suffice; for a million concurrent tasks, cache entries or connections, the
potential interactions between entities are $O(N^{2})$ and the complexity is infeasible, so
pooling, indexing, layering, asynchrony, batching and cache locality are forced in. That set
of techniques is exactly the full inventory of ``kernel technology''.

\begin{example}\label{ex:pg-memctx}
A relational database engine builds its own memory-context system because a
general-purpose allocator fragments and cannot track leaks under a load that ``creates and
destroys thousands of tuples per query''. It therefore constructs allocation and bulk
reclamation over a tree of contexts, which is an object pool inside the database domain. An
asynchronous runtime builds its own work-stealing scheduler because operating-system threads
are too heavy in the face of a million concurrent operations, and therefore constructs
$M{:}N$ scheduling, which is user-space time-sharing. A jump in order of magnitude forces
kernel technology in: this is the first necessity.
\end{example}

\section{The Six Functions}\label{sec:six}

\begin{definition}[The six functions]\label{def:six}
For a participant--resource structure $\Pi$, the \term{six functions} $\mathcal{K}_6$ are
defined as
\begin{enumerate}[label=(F\arabic*)]
\item \term{Identity and lifecycle}: the identity, creation, liveness and termination
      relations of entities;
\item \term{Allocation and reclamation}: the allocation and reclamation policies for finite
      resources;
\item \term{Scheduling}: the selection function that chooses the next executor among ready
      entities;
\item \term{Naming and addressing}: the mechanisms that make entities discoverable,
      referenceable and traversable;
\item \term{Communication and synchronisation}: the interaction primitives between entities
      and the composition rules for concurrent transitions;
\item \term{Fault isolation and recovery}: the containment relations of faults and the
      replayability of state.
\end{enumerate}
\end{definition}

\begin{theorem}[Conditional necessity of the six functions]\label{th:six}
Let $\Pi = (E, \mathcal{R}, \to, \mathrm{Ready}, \mathrm{Fault})$ be a
participant--resource structure whose operational semantics is well defined (that is, in
every reachable configuration the next transition is determined, or the non-determinism it
permits is explicitly specified). If $|E| > 1$ and $\mathcal{R}$ is finite and contended,
then the six functions of $\mathcal{K}_6$ have necessarily already been defined, either in
the system's own code or in its host (operating system, language runtime, framework).
\end{theorem}

\begin{proof}[Structural argument]
It is shown item by item that the absence of any single function renders the semantics
incomplete (an undefined behaviour exists).

\emph{(F1) Identity and lifecycle}: if entities were indistinguishable, $\mathrm{Ready}$
could not be defined as a function on the set of entities; if entities had no creation and
termination relations, the variation of $E$ over time would be unspecified and $\to$ would
be undefined at steps where ``an entity appears or disappears''.

\emph{(F2) Allocation and reclamation}: $\mathcal{R}$ is finite and contended, so ``more
resources are requested at once than can be supplied'' is a reachable state. Without an
allocation policy, $\to$ is undefined in that state. If the resources of terminated entities
are not reclaimed, available resources decrease monotonically and the system reaches
resource exhaustion in finitely many steps; $\to$ is then undefined in the forward-looking
sense.

\emph{(F3) Scheduling}: if $|E| > 1$ and several entities are ready at once, $\to$ is no
longer a function but a relation. Without a specified choice among ready entities, the
``permitted non-determinism'' is not delimited and the semantics is incomplete.

\emph{(F4) Naming and addressing}: an interaction primitive (see F5) must designate the
object of the interaction. Without a referencing mechanism, the object of an interaction
cannot be expressed.

\emph{(F5) Communication and synchronisation}: if entities interact, the global transition
is not the simple product of local transitions; without composition rules, the semantics of
concurrent transitions is undefined.

\emph{(F6) Fault isolation and recovery}: $\mathrm{Fault} \ne \varnothing$ (faults are a
physical and engineering given). Without containment relations, a local deviation has global
semantics, which is equivalent to the system having no well-defined behaviour under fault;
without recovery rules, the system never returns to a specified state after the first fault.

In sum, the absence of any one function destroys the premise of well-definedness. \qed
\end{proof}

\begin{remark}[The nature of the theorem]\label{rmk:six-nature}
Theorem \ref{th:six} is a \emph{conditional necessity} result, not an unconditional theorem.
Its conditions are $|E| > 1$, $\mathcal{R}$ finite and contended, and faults reachable. Its
proof is a structural enumeration argument, not a derivation in some formal system. It can
be falsified by exhibiting a well-defined system that satisfies the premises and lacks one
of the six. That judgement is itself open to challenge.
\end{remark}

The wording of Theorem \ref{th:six}, ``either in the system's own code or in its host'', is
essential. It separates two questions: whether the functions necessarily appear, and whether
they must be built in-house.

\section{Build or Adopt}\label{sec:buildadopt}

\begin{corollary}[Build-or-adopt dichotomy]\label{cor:buildadopt}
By Theorem \ref{th:six} the six functions must be defined, but where they are defined is a
design choice. That choice is decided by the \emph{domain criticality} of the invariants:
\begin{enumerate}[label=(\alph*)]
\item if the invariant carried by a function is \emph{domain-critical}, that is, the
      invariant can hold only under a guarantee owned by this system, the function must be
      built in-house;
\item if the invariant carried by a function is \emph{general}, that is, a general-purpose
      host can provide a guarantee of sufficient strength, the function should be adopted;
\item implementing a general function in-house produces the failure mode this paper calls
      \emph{reinvention of commodity parts}: a worse general-purpose mechanism
      re-implemented at a higher long-term cost of comprehension and maintenance.
\end{enumerate}
\end{corollary}

This criterion explains why the boundary between ``application'' and ``platform'' falls
exactly where it does: \emph{the set of invariants you are forced to own is the boundary of
your kernel}.

\begin{example}\label{ex:buildadopt}
A game engine must build its own entity--component storage and frame scheduling, because a
general-purpose host cannot provide the guarantee of a ``sub-millisecond deterministic frame
budget''. A database must build its own storage engine and write-ahead log, because a
general-purpose file system cannot provide the guarantee of ``recovery to a deterministic
state after a crash''. Conversely, a web application should adopt garbage collection, thread
scheduling and memory allocation, because the invariants of these functions are general and
building them in-house yields only a worse version.
\end{example}

\begin{remark}[A sharper criterion: monotonicity]\label{rmk:calm}
Corollary \ref{cor:buildadopt} decides build versus adopt by whether the invariant is
domain-critical. For functions of the ``coordination'' kind there is a sharper criterion that
can be checked statically: \emph{whether the computation is monotone}. Where monotonicity
holds \cite{calm2011} (outputs grow with inputs and do not depend on negation or on
aggregate global knowledge), coordination is not required and can be pushed down to the
participants; non-monotonicity (uniqueness constraints, negation, threshold aggregation,
non-commutative updates) forces some site to carry coordination.

This sharpens the claim for function (F5) (communication and synchronisation) and part of
(F3) (scheduling) of Section \ref{sec:six} from ``centralise once the scale is reached'' to
``centralise once non-monotonicity appears''. For monotone subsystems the $k$ of the
integration cost theorem (Theorem \ref{th:integration}) can be reduced to nearly a constant,
because participants need not agree on shared state.
\end{remark}

\section{Cross-Domain Comparison}\label{sec:instances}

Theorem \ref{th:six} predicts the pattern of Table \ref{tab:organs}: each column is a system
from a different company, a different language, a different era and a different domain, while
each row is the same organ. The table does not claim that these mechanisms can be identified
one with another; it compares them along the same dimension of \emph{management duty}.

\begin{table}[htbp]
\centering
\scriptsize
\caption{Independent implementations of the same set of management duties in different domains}\label{tab:organs}
\begin{tabular}{L{1.5cm}L{2.4cm}L{2.4cm}L{2.2cm}L{2.1cm}L{2.3cm}L{2.4cm}}
\toprule
\textbf{Function} & \textbf{Linux} & \textbf{Chromium} & \textbf{Tokio} & \textbf{Bevy}
& \textbf{Unreal} & \textbf{PostgreSQL} \\
\midrule
Identity/\newline lifecycle & task\_struct, inode, file & Frame, Document, Node &
Task, JoinHandle & Entity, Component & UObject, Actor & backend, xid, tuple \\
Allocation/\newline reclamation & buddy, slab, page cache & PartitionAlloc, Oilpan &
slab-style task allocation & archetype column store & FMemory, GC &
MemoryContext, shared\_buffers \\
Scheduling & EEVDF, workqueue & task queue, RendererScheduler &
work-stealing, waker & Schedule/System parallel executor & TaskGraph, GameThread &
postmaster, autovacuum \\
Naming/\newline addressing & kobject, sysfs, fd & DOM, Web IDL & Handle, waker &
Entity ID & FName, reflection metadata & OID, system catalog \\
Communi\-cation/\newline synchroni\-sation & syscall, netlink, IPC & Mojo IPC &
mpsc, oneshot &
Event, Resource & delegates, replication & libpq, LWLock, extension hooks \\
Isolation/\newline recovery & namespace, cgroup, OOM & site isolation, renderer process sandbox &
panic propagation, cancellation & state machine isolation & crash reporter &
MVCC snapshots, WAL replay \\
\bottomrule
\end{tabular}
\end{table}

Conversely, the \emph{form} these six organs take in different hosts varies widely, and that
variation is the concrete content of ``the trunk converges while the branches diverge'' in
Chapter \ref{ch:convergence}. The sources of that variation in form can be listed
systematically:

\begin{enumerate}[label=(\alph*)]
\item \term{Host memory model}. Systems with manual memory management (C/C++/Rust) must
      build their own allocator and even their own reclaimer (engine object reclamation,
      Blink's garbage collection in the browser); systems with managed memory spend the same
      effort on pause behaviour, because allocation is given and the pain point becomes the
      tail of the latency distribution.
\item \term{Host concurrency model}. The runtime model determines the shape of the
      scheduling organ: a game engine uses frame-driven cooperative scheduling (predictable
      latency over throughput), a server event-driven preemptive scheduling (throughput
      first). Frame schedulers and work-stealing schedulers are therefore two subspecies of
      one organ.
\item \term{Granularity of domain primitives}. A general-purpose kernel offers pages,
      threads, file descriptors and event polling; but a game engine needs frames,
      components and assets, a database needs buffer pools, transactions and version chains,
      and a browser needs document nodes, security contexts and cross-process objects. A
      mismatch in primitive granularity forces each domain to build its own infrastructure.
\item \term{Evolutionary stage}. These internal operating systems are themselves
      continuously refactored, and each refactoring converges towards ``a more complete set
      of the six functions'' rather than diverging from it.
\end{enumerate}

Even the operating system itself is circumvented by its users: high-performance network and
storage stacks redo I/O by polling in user space, asynchronous I/O interfaces are redesigned,
and the operating system cannot escape being ``reinvented'' by its own user space. This is
not anomalous; it is a direct consequence of Corollary \ref{cor:threshold}: when the
$n \times m$ of user space crosses the threshold, user space grows a new kernel.

\section{Transport and Re-Derivation}\label{sec:transport}

Table \ref{tab:organs} shows the same set of functions recurring across domains. But
``recurrence'' admits two entirely different explanations, with entirely different
predictions:

\begin{enumerate}[label=(\alph*)]
\item \term{Isomorphism}: the six functions are six implementations of one structure and can
      be translated into one another;
\item \term{Transport along the base}: the six functions are the fibres of some more general
      structure over different \emph{problem shapes}, and when the shape changes the
      structure must be re-derived, not carried over.
\end{enumerate}

The original treatment gives only (a). Distinguishing the two requires making the base
explicit, namely ``problem shape''.

\begin{definition}[Problem shape and level family]\label{def:fibration}
A pair $\pi : \mathcal{L} \to \mathcal{P}$ is a \term{level family}, where
\begin{itemize}
\item $\mathcal{P}$ is a set of \term{problem shapes}, determined by the resource form,
      concurrency structure, failure modes and time scales of the domain;
\item $\mathcal{L}$ is a set of layer implementations;
\item the fibre $\pi^{-1}(p)$ is the family of layer implementations feasible at problem
      shape $p$.
\end{itemize}
$\pi$ maps ``an implementation'' to ``what shape of problem it solves''. Reindexing $\pi$
gives \term{transport along the base}: given $p \to p'$, an implementation at $p$ is moved
to $p'$.
\end{definition}

\begin{proposition}[What transport preserves]\label{prop:transport}
Let $p \to p'$ be a map between problem shapes. If the map is an \emph{equivalence} (the same
resource form, concurrency structure, failure modes and time scales), then the \emph{slots}
of $\mathcal{K}_6$ (the question each function must answer) are preserved under transport;
the form of the implementation may change, but the slots are neither added nor removed. If
the map is not an equivalence, the semantic requirement of at least one slot changes, and
the implementation must be re-derived.
\end{proposition}

\begin{proof}[Argument]
Each item of $\mathcal{K}_6$ is uniquely determined by the semantic obligation of
Theorem \ref{th:six}: identity by the distinguishability of entities, allocation by the
finiteness and contention of resources, scheduling by the cardinality of the ready set,
naming by the object of interaction, communication by the mode of composition of concurrent
transitions, and isolation by the propagation range of faults. These six obligations depend
only on the four components of problem shape. When the map preserves the four components,
the obligation set is unchanged; when any component changes, the semantic requirement of the
corresponding obligation changes with it.
\end{proof}

\begin{corollary}[An observable discriminator]\label{cor:transport-test}
By Proposition \ref{prop:transport}, differences in form fall into two classes:
\emph{implementation differences of the same slot over different bases} (for instance
frame-driven and event-driven schedulers, whose bases have different time scales), and
\emph{differences of slot}. The former is not a counterexample; the latter is.
\end{corollary}

Corollary \ref{cor:transport-test} makes the middle proposition of Chapter
\ref{ch:convergence} (Hypothesis \ref{hyp:medium}) testable more rigorously: two systems
must first be aligned on the base, and only then compared slot by slot. Comparing modules
directly without aligning the base is the most common error in existing comparative studies
of architecture.

\begin{remark}[The status of this section]\label{rmk:fibration-status}
Definition \ref{def:fibration} and Proposition \ref{prop:transport} are a \emph{structural
proposal}, not an application of an established theorem of category theory. Their role is to
make explicit a question that has not previously been posed, whether the recurrence is
isomorphism or transport, and to give a discriminating corollary
(Corollary \ref{cor:transport-test}). Promoting $\pi$ to a strict fibration and proving the
properties of the reindexing functor is unfinished work (item O in
Chapter \ref{ch:objections}).
\end{remark}

\section{The Coarse-Graining Criterion}\label{sec:coarsegrain}

Theorem \ref{th:six} establishes that the six functions necessarily \emph{exist}, but not
that they are the most \emph{salient} structure. A stronger proposition is that the six
functions are the \emph{surviving set} of system structure under coarse-graining.

\begin{definition}[Coarse-graining sequence]\label{def:coarsegrain}
Let the architecture relational graph of a system be $\mathcal{G}$
(Definition \ref{def:relgraph}). A \term{coarse-graining} is a quotient map
$q : \mathcal{G} \to q(\mathcal{G})$ that merges components while preserving relation types.
A sequence of quotient maps $q_1, q_2, \dots$ (each step strictly coarser than the last) is
a \term{coarse-graining sequence}. A structure $\mathcal{S}$ \term{survives} under that
sequence if there exists $N$ such that for all $n \ge N$, $\mathcal{S}$ remains recognisable
in $q_n(\mathcal{G})$.
\end{definition}

\begin{conjecture}[The six functions are the relevant structure]\label{conj:relevant}
Under the coarse-graining sequence of a normally biased large system, the structures that
survive are exactly the slots of $\mathcal{K}_6$ and their relational graph; the remaining
structures (concrete modules, algorithm choices, data-layout details) are washed out at a
sufficiently coarse granularity.
\end{conjecture}

The value of Conjecture \ref{conj:relevant} is that it yields a falsification procedure
\emph{inside a single project}: no cross-project statistical comparison is needed, only
repeated coarse-graining of one code base, observing the surviving set. If the surviving set
is not the six functions, the conjecture is refuted and the convergence thesis of
Chapter \ref{ch:convergence} requires a different explanation.

\begin{principle}[Form of implementation]\label{pr:coarsegrain-impl}
Coarse-graining can be implemented entirely mechanically: starting from a module dependency
graph, three operations are applied repeatedly: (i) merge nodes whose fan-in and fan-out
fall below a threshold; (ii) collapse transitive dependency edges; (iii) erase nodes that
carry no invariant ($\mathrm{S}(\iota) = \varnothing$ in Definition \ref{def:enforce}).
Each step records the change in the surviving structure. The whole experiment needs one
repository and one dependency graph, with no labels and no manual classification, and is
therefore unaffected by the ``freedom in choosing a duty-classification system'' noted in
Section \ref{sec:limits}.
\end{principle}

Chapter \ref{ch:program} states Conjecture \ref{conj:relevant} as Hypothesis $H_{10}$.

\begin{remark}[Relation to renormalisation]\label{rmk:rg}
The operations of Principle \ref{pr:coarsegrain-impl} are formally the same as
coarse-graining in the renormalisation group: short-range degrees of freedom are eliminated
and one observes whether long-range structure converges to a fixed point. The fixed point
here is the six-function slots. It must be stated plainly that this does \emph{not} mean
that software possesses a renormalisation-group flow: software has no separation of scales
(the discussion preceding Principle \ref{pr:vm}), hence no exponential decay of genuinely
irrelevant operators. What is the same is the \emph{operation}, not the \emph{mechanism}.
\end{remark}

\section{Virtualisation and Deep Modules}\label{sec:virtualisation}

The phrase ``compress complexity and provide the upper layers with a high-level language and
symbols'' can be made precise. The core work of every infrastructure layer is one and the
same action:

\begin{principle}[Virtualisation]\label{pr:vm}
An infrastructure layer translates what is \emph{physical, continuous and volatile} into
\emph{discrete, stable, composable} domain-level symbols.
\end{principle}

\begin{example}\label{ex:vm}
Operating system: physical memory $\to$ virtual address space; block devices $\to$ files;
hardware interrupts $\to$ signals.
Database: physical pages $\to$ multi-version tuples $\to$ ``rows'' in SQL.
Browser: byte streams $\to$ document trees; socket events $\to$ event-loop callbacks.
Game engine: raw bytes in memory $\to$ components and assets; the time of one frame $\to$ a
schedule table.
Asynchronous runtime: event-poll bitmap $\to$ the Future of \code{async}/\code{await}.
\end{example}

In the vocabulary of module design \cite{parnas1972,ousterhout2018}, these are all \emph{deep
modules}: minimal interface, very deep implementation. Behind one allocation call stands an
entire allocator system; behind one task creation stands an entire scheduler. The duty of
an infrastructure layer is to keep producing deep modules, sealing complexity below and
exposing clean symbols above.

This connects directly to the theory of comprehension in Chapter \ref{ch:comprehension}: a
deep module is precisely a level of very high compression ratio
(Definition \ref{def:compression}) that quotients away a large amount of lower-level
difference while retaining the predictive power the upper layer needs. It also explains why
good infrastructure is at once ``high-performance'' and ``easy to understand'': both are
consequences of the same quotient map.

\section{Functional Differentiation in Agent Systems}\label{sec:agent}

If Theorem \ref{th:six} holds, the most recently emerged complex computational domain should
be re-enacting this evolution. Table \ref{tab:agent} gives the observation.

\begin{table}[htbp]
\centering
\small
\caption{Agent systems rediscovering the six functions}\label{tab:agent}
\begin{tabular}{L{3.4cm}L{11.0cm}}
\toprule
\textbf{Function} & \textbf{Counterpart in the agent domain} \\
\midrule
Allocation and reclamation & the context window as fast storage; retrieval augmentation as
page-fault handling; external memory as swap space; context compression as memory
compaction \\
Communication and synchronisation & tool calls or function calls as system-call
primitives \\
Scheduling & the agent main loop, multi-agent orchestration \\
Identity and lifecycle & lifecycle management of sessions, tasks and subgoals \\
Naming and addressing & tool registration and discovery protocols \\
Isolation and recovery & sandboxing, retries, checkpoints and rollback \\
\bottomrule
\end{tabular}
\end{table}

The value of this observation is that it is a \emph{natural experiment}: no one is
deliberately imitating a kernel; complexity itself has pressed out the same anatomy. It also
forms the empirical basis of Hypothesis $H_3$ (Chapter \ref{ch:program}).

Theorems \ref{th:integration} and \ref{th:six} hold under their stated premises; what can be
challenged is the premises and the rigour of the structural argument. What is empirically
falsifiable is whether the appearance of the six functions is predicted by $n \times m$
rather than by lines of code (Hypotheses $H_3$ and $H_{10}$), and whether there exists a
well-defined system that satisfies the premises yet lacks one of the functions over the long
run.
